% Part: history
% Chapter: biographies
% Section: gerhard-gentzen

\documentclass[../../../include/open-logic-section]{subfiles}

\begin{document}

\olfileid[pt]{his}{bio}{gen}

\olsection{Gerhard Gentzen}

\olphoto{gentzen-gerhard}{Gerhard Gentzen}

Gerhard Gentzen é conhecido principalmente como o criador da teoria estrutural
da prova e, especificamente, pela criação dos sistemas de !!{derivation} de
dedução natural e de cálculo de sequentes. Ele nasceu em 24 de novembro de 1909
em Greifswald, na Alemanha. Gerhard foi educado em casa por três anos antes de
frequentar o ensino preparatório, onde estava atrás da maioria de seus colegas
em termos de instrução. Apesar disso, era um aluno brilhante e demonstrava
forte aptidão para a matemática. Seus interesses eram variados e ele, por
exemplo, também escrevia poemas para a mãe e peças para o teatro da escola.

Gentzen começou seus estudos universitários na Universidade de Greifswald,
mas passou por G\"{o}ttingen, Munique e Berlim. Recebeu seu doutorado em 1933
pela Universidade de G\"{o}ttingen, sob a orientação de Hermann Weyl. (Paul
Bernays supervisionou a maior parte de seu trabalho, mas foi demitido da
universidade pelos nazistas.) Em 1934, Gentzen começou a trabalhar como
assistente de David Hilbert. Nesse mesmo ano, desenvolveu os sistemas de
!!{derivation} de cálculo de sequentes e de dedução natural, em seus artigos
\emph{Untersuchungen \"{u}ber das logische Schlie\ss en I--II
  [Investigações sobre a Dedução Lógica I--II]}. Ele provou a
consistência dos axiomas de Peano em 1936.

A relação de Gentzen com os nazistas é complicada. Ao mesmo tempo em que seu
mentor Bernays foi forçado a deixar a Alemanha, Gentzen ingressou no braço
universitário da SA, a organização paramilitar nazista. Como muitos alemães,
foi membro do partido nazista. Durante a guerra, serviu como oficial de
telecomunicações na unidade de inteligência aérea. No entanto, em 1942 foi
dispensado do serviço devido a um colapso nervoso. Não está claro se as
lealdades de Gentzen estavam com o partido nazista ou se ele ingressou no
partido a fim de garantir sucesso acadêmico.

Em 1943, foi oferecida a Gentzen uma posição acadêmica no Instituto
Matemático da Universidade Alemã de Praga, que ele aceitou. No entanto, em
1945 os cidadãos de Praga se revoltaram contra a ocupação alemã. As forças
soviéticas chegaram à cidade e prenderam todos os professores da universidade.
Por causa de sua filiação a organizações nazistas, Gentzen foi levado a um
campo de trabalhos forçados. Ele morreu de desnutrição em sua cela em 4 de
agosto de 1945, aos 35 anos.

\begin{reading}
Para uma biografia completa de Gentzen, veja \citet{Menzler-Trott2007}. Uma
leitura interessante sobre matemáticos sob o regime nazista, que traz uma
breve nota sobre a vida de Gentzen, é dada por \citet{Segal2014}.
Os artigos de Gentzen sobre dedução lógica estão disponíveis no alemão
original \citep{Gentzen1935a,Gentzen1935b}. Traduções para o inglês dos
artigos de Gentzen foram reunidas em um único volume por
\citet{Gentzen1969}, que também inclui um esboço biográfico.
\end{reading}

\end{document}
